{\Large {\bfseries Е}лектро{\bfseries м}агнітне {\bfseries в}ипромінювання (EMR)}
\begin{itemize}
\item EMR випромінюється дискретними одиницями, що називаються фотонами, але має властивості хвиль. EMR може бути створене коливанням або прискоренням електричного заряду чи магнітного поля. EMR поширюється у просторі зі швидкістю світла ($c=2.997 924 58 \times 10^{8}$ \si{\metre\per\second} ). EMR складається з електричного та магнітного полів, що коливаються перпендикулярно одне до одного з певною довжиною хвилі.

\rput[t]{0}(0.7,-.5){\input{pictures/emr.tex}}
\rput[t]{0}(2.85,-.5){\input{pictures/emr_particle.tex}}

\vspace{1.45in}

\item Корпускулярна природа проявляється, коли слабко освітлений сонячний елемент випускає окремі електрони. Хвильова природа демонструється експериментом з подвійною щілиною, що показує погашення та додавання хвиль.

% \item Much of the EMR properties are based on theories since we can only see the effects of EMR and not the actual photon or wave itself.

% \item Einstein theorized that the speed of light is the fastest that anything can travel. He has not been proven wrong.

\item Довжина хвилі EMR змінюється, коли джерело віддаляється або наближається. Червоне зміщення робить зірки та галактики, що швидко віддаляються, більш червоними.

% \item We only have full electronic control over frequencies in the microwave range and lower. Higher frequencies must be created by waiting for the energy to be released from elements as photons. We can either pump energy into the elements (ex. heating a rock with visible EMR and letting it release infrared EMR) or let it naturally escape (ex. uranium decay).
%
% \item We can only see the visible spectrum. All other bands of the spectrum are depicted as hatched colors \psframebox[fillstyle=none,linestyle=none,framesep=0in]{\psframe[linearc=0,framearc=0,fillstyle=crosshatch,linewidth=0pt,linestyle=none, hatchwidth=2pt, hatchsep=1.5pt,hatchcolor=white](0,-.04)(.4,.1)}\hspace{0.4in}.


\end{itemize}
